\chapter{Тёплиц-конвейер свежих вспомогательных систем полного шумового процесса}
\label{ch:v8-full-noise-toeplitz-ancilla-chain-dilation-gate}

Ранний оператор числа и двусторонний сдвиг
\begin{equation}
 Ne_n=ne_n,\qquad U^ke_n=e_{n+k},\qquad [N,U^k]=kU^k
\end{equation}
дают абсолютный целый адрес такта. Используем этот каркас для полной
столкновительной среды
\begin{equation}
 \mathcal K_{\rm cell}=\mathbb C^{43}
 =\mathbb C|0\rangle\oplus\mathbb C^{42}_{\rm jump}
\end{equation}
и двусторонней цепи относительно опорного вакуумного состояния
\begin{equation}
 \mathcal K_{\rm chain}
 =\bigotimes_{m\in\mathbb Z}(\mathcal K_{\rm cell},|0\rangle).
\end{equation}

Пусть $U_{\rm col}^{(0)}$ есть унитарий звёздного взаимодействия системы с
нулевой ячейкой, а $S_{\rm chain}$ сдвигает ячейки в сторону системы.
Один фиксированный глобальный такт равен
\begin{equation}
 V=(I_{21}\otimes S_{\rm chain})U_{\rm col}^{(0)}.
\end{equation}
После взаимодействия использованная ячейка уходит, а следующий такт
подводит ранее не затронутую вакуумную ячейку. Для любого конечного $n$ это
утверждение живёт на конечном цилиндре из $n$ ячеек и не требует вычислять
бесконечную матрицу.

\begin{theorem}[Точное автономное повторение канала]
Для тензорного вакуума цепи и любого конечного $n\geq0$
\begin{equation}
 \Tr_{\rm chain}\!\left[
 V^n(\rho_0\otimes\omega_{\rm vac})V^{-n}\right]
 =\Phi_h^n(\rho_0).
\end{equation}
Один и тот же унитарий $V$ реализует все такты; внешняя операция сброс или
подмена вспомогательных систем между шагами не требуется.
\end{theorem}

Калибровочное действие продолжается одинаково на каждую 43-мерную ячейку. Поскольку
42-мерный репер скачков замкнут под калибровочными коммутаторами, локальное столкновение и
сдвиг одинаковых ячеек дают ковариантный глобальный такт. При слабом
масштабировании сохраняется прежний предел $\Phi_{u/n}^n\to e^{u\mathcal
L_{42}}$.

\begin{nogocriterion}[Граница сильной автономности]
Конструкция автономна как дискретный Флоке-процесс, но использует заранее
приготовленную бесконечную цепь в тензорном вакууме. Её состояние, локальная
энергия, происхождение сдвига из одного стационарного гамильтониан и
длительность такта ещё не выведены родительского действия. Поэтому результат нельзя
называть полностью самопорождающимся физическим резервуаром.
\end{nogocriterion}

\begin{protocol}
\begin{enumerate}
 \item система: \textbf{$21$};
 \item метки скачков: \textbf{$42$};
 \item ячейка: \textbf{$43$};
 \item абсолютный счётчик: \textbf{$[N,U^k]=kU^k$};
 \item один глобальный такт: \textbf{унитарен};
 \item повторное использование старой ячейки: \textbf{нет};
 \item редуцированная динамика: \textbf{$\Phi_h^n$ точно};
 \item калибровочная ковариантность: \textbf{по ячейкам};
 \item внешний сброс каждого шага: \textbf{не нужен};
 \item вакуумная цепь из родительского действия: \textbf{не выведена};
 \item не зависящий от времени локальный гамильтониан: \textbf{не выведен};
 \item физическая длительность такта: \textbf{не выведена}.
\end{enumerate}
\end{protocol}